% ============================================================
% DOKUMENTKLASSE UND SPRACHE
% ============================================================
% für elektronische Veröffentlichung (Standard)
\documentclass[11pt,oneside,paper=A4,DIV=11,BCOR=0mm,abstract=true,headsepline,headings=normal,ngerman]{scrreprt}
% für Druck (doppelseitig)
%\documentclass[11pt,twoside,paper=A4,DIV=11,BCOR=8mm,abstract=true,headsepline,headings=normal,ngerman]{scrreprt}

\usepackage{babel}
\usepackage[utf8]{inputenc}

% ============================================================
% SCHRIFTART
% ============================================================
\usepackage{libertine}
\usepackage{libertinust1math}
\usepackage[T1]{fontenc}

% ============================================================
% TYPOGRAFIE
% ============================================================
\usepackage[auto]{microtype}
\clubpenalty = 10000
\widowpenalty = 10000
\displaywidowpenalty = 10000

% ============================================================
% MATHE, SYMBOLE, EINHEITEN, CHEMIE
% ============================================================
\usepackage{amsmath}
\usepackage{amsxtra}
\usepackage{eurosym}
\usepackage{siunitx}
\sisetup{locale=DE}
\usepackage[version=4]{mhchem}
\usepackage{chemfig}

% ============================================================
% LOGIK-HILFSPAKET (für nomencl-Bedingungen \ifstrequal)
% ============================================================
\usepackage{etoolbox} % NEU: wird von \ifstrequal in der Nomenklatur benötigt

% ============================================================
% FARBEN UND GRAFIKPROGRAMME
% ============================================================
\usepackage{graphicx}
\usepackage[table]{xcolor} % ANGEPASST: Option "table" ergänzt für \columncolor in Tabellen
\usepackage{tikz}
\usetikzlibrary{fit, backgrounds, positioning, decorations.pathreplacing, arrows.meta, positioning}
\usepackage[edges]{forest} % ANGEPASST: Duplikat entfernt, edges-Option direkt hier geladen
\usepackage{dirtree}

% ============================================================
% TABELLEN UND FLOATS
% ============================================================
\usepackage{multirow,multicol,booktabs}
\usepackage{threeparttable}
\usepackage{longtable}
\usepackage{xltabular}   % Seitenumbrüche + flexible Spalten
\usepackage{ltablex}
\usepackage{array}       % erweiterte Spaltendefinitionen
\usepackage{float}       % Positionierung (H-Option)
\usepackage{seqsplit}    % Umbruch langer Zeichenketten
\usepackage{rotating}
\usepackage{pdflscape}   % Drehen ganzer PDF-Seiten

% ============================================================
% ABBILDUNGEN, TEILABBILDUNGEN, EINGEBETTETE DOKUMENTE
% ============================================================
\usepackage{pdfpages}
\usepackage{subcaption} % ANGEPASST: ersetzt subfig, siehe Hinweis unten
\captionsetup[subtable]{position=top}

% ============================================================
% TEXTGESTALTUNG UND CODE
% ============================================================
\usepackage{listings}
\usepackage{minted}
\usepackage{blindtext}
\usepackage{csquotes}
\usepackage[normalem]{ulem} % durchgestrichener Text
\usepackage{url}
\urlstyle{same}
\usepackage{comment}

% ============================================================
% FUSSNOTEN
% ============================================================
\usepackage{footnote}
\makesavenoteenv{tabular}
\usepackage{chngcntr}
\counterwithout{footnote}{section} % entkoppelt Fußnoten vom Hauptkapitel

% Fußnotenmarkierung linksbündig statt hochgestellt
\deffootnote{1.5em}{1em}{%
    \makebox[1.5em][l]{\thefootnotemark}%
}
\interfootnotelinepenalty=10000 % Fußnoten nicht umbrechen

% ============================================================
% LITERATURVERZEICHNIS
% ============================================================
\usepackage[
    backend=biber,
    style=numeric,
    sorting=none,
    citestyle=numeric-comp,
    maxbibnames=10,
    giveninits=true
]{biblatex}

\renewcommand*{\labelnamepunct}{\addcolon\addspace}

% Speicherort der Literaturangaben (*.bib-Datei)
\bibliography{literatur/literaturdatenbank}

% ============================================================
% NOMENKLATUR (ABKÜRZUNGS- UND SYMBOLVERZEICHNIS)
% ============================================================
\usepackage[intoc, german]{nomencl}
\renewcommand{\nomname}{Abkürzungs- und Symbolverzeichnis}
\setlength{\nomitemsep}{-\parsep}
\newlength{\nomgroupstartsep}
\setlength{\nomgroupstartsep}{18pt}
\newcommand{\nomenclheader}[1]{\item[\hspace*{-\itemindent}\itshape#1]{\ } \vspace{6pt}\par}

\renewcommand\nomgroup[1]{
 \itemsep\nomgroupstartsep
  \ifstrequal{#1}{A}{\nomenclheader{Abkürzungen}}{
  \ifstrequal{#1}{E}{\nomenclheader{Lateinische Symbole}}{
  \ifstrequal{#1}{G}{\nomenclheader{Griechische Symbole}}{
  \ifstrequal{#1}{O}{\nomenclheader{Operatoren}}{
  \ifstrequal{#1}{X}{\nomenclheader{Tiefgestellte Indizes}}{
  \ifstrequal{#1}{Y}{\nomenclheader{Hochgestellte Indizes}}{}}}}}}
 \itemsep\nomitemsep
}

% ============================================================
% GESTALTUNG VON BILDUNTERSCHRIFTEN, TABELLEN, TITELSEITE
% ============================================================
\addtokomafont{caption}{\small}
\setkomafont{captionlabel}{\sffamily\bfseries}
\setkomafont{author}{\large}
\setkomafont{date}{\large}
\setkomafont{publishers}{\large}
\setkomafont{dedication}{\normalsize}
\addtokomafont{dedication}{\raggedright}

\renewcaptionname{ngerman}{\figurename}{Abb.}
\renewcaptionname{ngerman}{\tablename}{Tab.}

% Tabellenumgebungen mit Schriftgröße 10 und 7
\newenvironment{tabular10}{%
    \fontsize{10}{12}\selectfont\tabular
}{%
    \endtabular
}

\newenvironment{tabular7}{%
    \fontsize{7}{12}\selectfont\tabular
}{%
    \endtabular
}

% ============================================================
% EXTERNE DATENDATEI (Titel, Autor, etc.)
% ============================================================
%\def\VARtyp{Bachelorarbeit}
\def\VARtyp{Masterarbeit}

\def\VARtitel{Videoquellenerkennung für MP4-Container mittels Similarity Hashing}

\def\VARuntertitel{}

\def\VARname{Lucie Mechelk}

\def\VARmatrikelnummer{1254324}
\def\VARstudiengang{Cyber-Security}

\def\VARabgabedatum{\today}



% ============================================================
% VERWEISE, REFERENZEN, PDF-METADATEN
% ============================================================
% WICHTIG: hyperref muss möglichst spät geladen werden,
% cleveref direkt danach
\usepackage[
    pdftitle={\VARtitel},
    pdfsubject={\VARtyp},
    pdfauthor={\VARname},
    pdfkeywords={},
    hidelinks % Links nicht einrahmen
]{hyperref}

\usepackage[german]{cleveref}